\documentclass[a4paper,11pt]{ltjsarticle}
\usepackage[haranoaji]{luatexja-preset}
\usepackage{amsmath,amssymb}
\usepackage[top=12mm,bottom=10mm,left=14mm,right=14mm]{geometry}
\usepackage{array}
\usepackage{tikz}
\usetikzlibrary{calc}
\pagestyle{empty}
\setlength{\parindent}{0pt}
\newlength{\qcol}\setlength{\qcol}{0.47\textwidth}
\newlength{\qgap}
\newlength{\anscolw}\setlength{\anscolw}{0.30\textwidth}
\newlength{\ansh}
\newcommand{\Q}[2]{\parbox[t]{\qcol}{\raggedright (#1)\hspace{0.6em}$\displaystyle #2$}}
\newcommand{\QR}[4]{\Q{#1}{#2}\hfill\Q{#3}{#4}\par\vspace{\qgap}}
\newcommand{\anscell}[1]{\rule[-\ansdrop]{0pt}{\ansh}(#1)}
\newlength{\ansdrop}

\begin{document}
{\large\bfseries 計算問題　14}\hfill 名前(\underline{\hspace{32mm}})\par\vspace{3mm}
■（1）〜（16）の計算をしなさい。（17）〜（20）は方程式を解きなさい。\par\vspace{3mm}
\setlength{\qgap}{5.4mm}
\QR{1}{2(3x+1)-(4x-5)}{2}{(3a+1)+(-7a-6)}
\QR{3}{12x\div(-3)}{4}{(-5)\times13a}
\QR{5}{3-(-6)+(-2)^3}{6}{(-48)\div(-4)}
\QR{7}{-7+(-13)}{8}{(-11)\times(-9)}
\QR{9}{\dfrac{x-1}{3}\times9}{10}{(30x-9)\div\dfrac{3}{2}}
\QR{11}{1-3+5-7}{12}{5\div4\times(-16)}
\QR{13}{-4^2+9\div(-3)}{14}{\dfrac{1}{4}(4x-8)+\dfrac{1}{3}(6x+3)}
\QR{15}{(-5)-(-15)}{16}{-12\div\{-2\times(4-6)\}}
\QR{17}{-3(x+5)=-2(x-4)}{18}{7x-3=11}
\QR{19}{0.2x+0.8=3}{20}{\dfrac{5}{2}x+1=\dfrac{2}{3}x}
\vspace{1mm}
\Q{21}{75\times8}\par\vspace{3mm}
\setlength{\ansh}{9.5mm}\setlength{\ansdrop}{6mm}%
\begin{center}\begin{tabular}{|p{\anscolw}|p{\anscolw}|p{\anscolw}|}\hline
\anscell{1} & \anscell{2} & \anscell{3} \\\hline
\anscell{4} & \anscell{5} & \anscell{6} \\\hline
\anscell{7} & \anscell{8} & \anscell{9} \\\hline
\anscell{10} & \anscell{11} & \anscell{12} \\\hline
\anscell{13} & \anscell{14} & \anscell{15} \\\hline
\anscell{16} & \anscell{17} & \anscell{18} \\\hline
\anscell{19} & \anscell{20} & \anscell{21} \\\hline
\end{tabular}\end{center}

\end{document}
